% !PRINTED 682
\weberpage{682}
But by the multiplication theorem of determinants

\weq{3}{\begin{vmatrix} \lambda_{1,1}, & \cdots, & \lambda_{\nu-1,1} \\ \cdot & \cdot & \cdot \\ \lambda_{1,\nu-1}, & \cdots, & \lambda_{\nu-1,\nu-1} \end{vmatrix}
=
\begin{vmatrix} \xi_{1,1}, & \cdots, & \xi_{1,\nu-1} \\ \cdot & \cdot & \cdot \\ \xi_{\nu-1,1}, & \cdots, & \xi_{\nu-1,\nu-1} \end{vmatrix}
\begin{vmatrix} l_{1,1}, & \cdots, & l_{\nu-1,1} \\ \cdot & \cdot & \cdot \\ l_{1,\nu-1}, & \cdots, & l_{\nu-1,\nu-1} \end{vmatrix},}

or, if we set

\[
\Lambda_{\nu-1} = \Sigma \pm \lambda_{1,1}\lambda_{2,2} \cdots \lambda_{\nu-1,\nu-1}
\]

and observe that the determinant

\[
\Sigma \pm \xi_{1,1}\xi_{2,2} \cdots \xi_{\nu-1,\nu-1}
\]

is a rational number with the denominator $m^{\nu-1}$:

\weq{4}{\Lambda_{\nu-1} = \frac{a}{m^{\nu-1}} \, L(\varepsilon_1, \varepsilon_2, \ldots, \varepsilon_{\nu-1}),}

in which $a$ is an integral rational number.

From this it follows that the newly introduced units form an
independent system always and only then, when the
determinant of the $\xi_{i,k}$, that is, the number $a$, is different from zero.

Now among a finite or infinite number of
non-vanishing rational fractions with the same denominator
there is always a least in absolute value, and consequently
we can, by (4), choose the new system of independent units
so that its regulator $\pm \Lambda_{\nu-1}$ becomes \gesperrt{as small as possible}.
Such a system we call a \gesperrt{fundamental system
of units}, and the minimal value of the regulator itself,
that is, the regulator of a fundamental system, we call the
\gesperrt{regulator of the field}. We now assume that our
system $\varepsilon_1, \varepsilon_2, \ldots, \varepsilon_{\nu-1}$ is itself a fundamental system.

Under this assumption it can be proved that
all exponents of units must be \gesperrt{integers}.

For if we assume that there exists a unit $\varepsilon$
whose exponents, determined by the formulae \S.\,186, (5), are not
all integers, then by \S.\,186, 4) there is also a
unit $\varepsilon_0$ whose exponents are proper fractions which do not all
vanish.

If by $\lambda_1, \lambda_2, \ldots, \lambda_{\nu-1}$ in the formulae (5),
\S.\,186 we understand the system of conjugate logarithms of $\varepsilon_0$, then there
 results

\[
L(\varepsilon_0, \varepsilon_2, \ldots, \varepsilon_{\nu-1}) = \xi_1\, L(\varepsilon_1, \varepsilon_2, \ldots, \varepsilon_{\nu-1}),
\]
%CONT>>
